\begin{abstract}
Η ραγδαία εξάπλωση των συσκευών του Διαδικτύου των Πραγμάτων (IoT) έχει καταστήσει τη σύσταση αυτοματισμών (κανόνων της μορφής «αν συμβεί το Α, τότε εκτέλεσε το Β») ένα ολοένα και πιο σημαντικό πρόβλημα. Τα δεδομένα που απαιτούνται για την εκπαίδευση ενός τέτοιου συστήματος είναι όμως εξ ορισμού ευαίσθητα, γεγονός που καθιστά την κεντρικοποιημένη συλλογή τους προβληματική. Η παρούσα εργασία μελετά το πρόβλημα της σύστασης κανόνων IoT σε ομοσπονδιακό πλαίσιο, ορμώμενη από τον ανοικτό διαγωνισμό της εταιρείας Wyze και τη λύση της δεύτερης θέσης.

Το πρόβλημα επαναδιατυπώνεται ως ομοσπονδιακή πρόβλεψη συνδέσεων σε γράφους ανά χρήστη και αναπτύσσονται τρεις διαδοχικές εκδόσεις μοντέλου αυξανόμενης πολυπλοκότητας, καταλήγοντας σε ένα relational μοντέλο που συνδυάζει τον CompGCN encoder, τον ComplEx decoder και τη συνάρτηση σφάλματος InfoNCE. Το μοντέλο αυτό βελτιώνει την κεντρικοποιημένη λύση αναφοράς κατά 71\% στη μετρική MRR. Η συστηματική αξιολόγηση σε τέσσερα σενάρια ενορχήστρωσης (κεντρικοποιημένο, cross-silo και δύο παραλλαγές cross-device) αποκαλύπτει ότι η βέλτιστη επιλογή μοντέλου εξαρτάται από τον τύπο της ενορχήστρωσης: το ισχυρότερο μοντέλο κυριαρχεί παντού, εκτός από το πιο κατανεμημένο σενάριο, όπου μια απλούστερη έκδοση αποδεικνύεται ανθεκτικότερη.

\vspace{1em}
\noindent\textbf{Λέξεις-κλειδιά:} Ομοσπονδιακή Μάθηση, Νευρωνικά Δίκτυα Γράφων, Συστήματα Σύστασης, Σύσταση Κανόνων, Διαδίκτυο των Πραγμάτων, Πρόβλεψη Συνδέσεων, Μη-IID Δεδομένα
\end{abstract}
